\chapter{Родитель спектрального полюса трансмутированной моды}
\label{ch:v10-discrete-resolution-transmuted-spectral-pole}

Пусть в клеточных единицах
\begin{equation}
 k_*\ell_{\rm cell}=e^{-32\pi^2/3},\qquad
 m_*^2\ell_{\rm cell}^2=e^{-64\pi^2/3}.
\end{equation}
Условный минковский обратный пропагатор
\begin{equation}
 \Gamma_*^{(2)}(q)=q-e^{-64\pi^2/3}
\end{equation}
имеет простой нуль при $q=m_*^2\ell_{
m cell}^2$, производную единицу и
положительный единичный вычет. Двумерный спектральный свидетель
$\operatorname{diag}(e^{-64\pi^2/3},1)$ имеет ранговый проектор на эту
моду. Строгий родитель коэффициента массы имеет гессиан $(1)$.

Это закрывает архитектуру полюса, но не его происхождение. Унаследованный
обратный пропагатор остаётся безмассовым, а положительный $b=2$ создаёт
ультрафиолетовый полюс Ландау на множителе $e^{+32\pi^2/3}$, а не
инфракрасный массовый нуль на $e^{-32\pi^2/3}$. Переход к обратной
экспоненте требует нового массового члена самоэнергии.

\begin{theorem}[Условный здоровый полюс существует]
Требуемое разрешение допускает простой спектральный полюс с положительным
вычетом и однозначным проектором. Но этот полюс не следует из
унаследованной положительной beta-функции.
\end{theorem}

\begin{nogocriterion}[UV-полюс Ландау не является IR-массой]
Обратная экспонента и нуль пропагатора не возникают из факта расходимости
связи в ультрафиолете. Необходим независимый механизм генерации массового
члена или инфракрасной самоэнергии.
\end{nogocriterion}

\begin{protocol}\begin{enumerate}
\item условная архитектура полюса: \textbf{$10/10$};
\item простой полюс и положительный вычет: \textbf{$6/6$};
\item унаследованный IR-массовый член: \textbf{$0/1$};
\item физическое происхождение полюса: \textbf{$0/3$};
\item следующий гейт: \textbf{аудит источников IR-массового члена}.
\end{enumerate}\end{protocol}

Машинный сертификат сохранён в
\path{s2t_v10_cell_birth_four_volume_nonequilibrium_bath_discrete_resolution_transmuted_mode_spectral_pole_parent_origin_gate_results.json}.